\documentclass[10pt]{article}
\usepackage[a4paper,margin=26mm]{geometry}
\usepackage{amsmath,amssymb,amsthm,hyperref}
\newtheorem{theorem}{Theorem}
\newtheorem{lemma}{Lemma}
\newcommand{\MA}{\operatorname{MA}}
\newcommand{\vol}{\operatorname{vol}}
\title{Weighted identification at finite divisorial outputs by local squeezing}
\author{An unreviewed conditional bridge for the original big-class equation}
\date{6 October 2026}
\begin{document}
\maketitle
\begin{abstract}
We propose a local squeezing argument identifying the mixed-moving
algebraic moment with the original relative integral when the output is
supported at finitely many divisorial points. The varying quotient of two
complete-system section metrics need not decrease or converge uniformly.
Two fixed-degree bounds, upper and lower semicontinuity near the output
atoms, and continuous localization suffice instead. We state exactly the
genuine envelope and output foundations required. Subject to those
foundations and the separate nef moment estimate, the identification
extends to every original relative finite-energy test for these atomic
outputs. General output pole-nullness, weighted convergence for non-atomic
outputs, actual arbitrary-target attainment and comparison remain open.
\end{abstract}

\paragraph{Fixed actual domain.}
Let $X$ be smooth integral projective of dimension $n\ge1$ over any
algebraically closed trivially valued characteristic-zero field.
Let $L$ be an actual big rational bundle, $K=X^{\rm an}$ and
$V=\vol(L)>0$. All full section systems and their common resolutions
are actual. Write $v_0$ for the generic trivial valuation and
$V_\theta$ for the normalized genuine reference.

This paper is conditional on the following exact foundations, which
must be reconstructed for original-scope acceptance:
\begin{enumerate}
\item The genuine numerical cone has decreasing closure and model
obstacle envelopes. For every rational bounded model obstacle $f_Z$,
$u_Z=P_\theta(f_Z)$ equals the usc regularization of its full complete
relative-system supremum. Its stages $\phi_{m,Z}$ increase in cofinal
divisibility order, are genuine and lie below $u_Z$ on the full carrier.
For $Z=0$, $u_0=V_\theta$. All are bounded above by a common constant.
\item Divisorial points are nonnegligible for a fixed actual ample
class $\omega$ containing $\theta$ in its psh cone. Consequently an
increasing uniformly upper-bounded family of these genuine potentials
has no usc-regularization gap at any divisorial point. This is an ample
background statement, together with its actual-class application;
no countability of $k$ is assumed.
\item Resolving the common generic complete system gives
$p_m^*L=A_m+T_m$, with $mA_m$ generated and $T_m\ge0$. Resolving the
relative ideals gives bounded nef model potentials $\psi_{m,Z}$ on
$A_m$ with
\[
 \phi_{m,Z}\circ p_m^{\rm an}=q_{T_m}+\psi_{m,Z},\qquad
 \phi_{m,0}\circ p_m^{\rm an}=q_{T_m}.
 \tag{1}
\]
For $\mu_m=(p_m^{\rm an})_*\MA_{A_m,\rm raw}(\psi_{m,D})$,
$\mu_m(K)=A_m^n\to V$ and $\mu_m\to\mu_D$ weakly, where
$\mu_D=\sum_{j=1}^N c_j\delta_{v_j}$ is the genuine
\emph{target-independent} model-envelope output, $c_j\ge0$,
$\sum c_j=V$, and all $v_j$ are divisorial points of the original
function field. Every auxiliary model measure is also supported at
such valuations after pushforward, with its actual central multiplicities.
\item The genuine raw energy is concave, constant-equivariant and
continuous on decreasing nets. Decreasing model-majorant envelopes are
cofinal for every genuine finite-energy potential. The shared full-system
raw model energies and generic values tend to their envelope energy and
generic value. The nef-model moment estimate and mixed-moving limit
stated below have been independently reconstructed.
\end{enumerate}
No output of a general unbounded potential, no relative-pole-nullness
for such an output and no non-atomic weighted limit is among these
foundations. Finite atomicity is an essential hypothesis of the result.

\section{An abstract localized squeeze}
\begin{lemma}\label{squeeze}
Let $K$ be compact Hausdorff, let positive Radon measures $\nu_i\to\nu$
weakly with converging masses, and suppose $\nu$ is supported by a finite
set $S$. Let functions $a_i$ be measurable on the supports of $\nu_i$,
with $|a_i|\le C$. For each degree $q$, suppose there is a neighborhood
$U_q$ of $S$ and bounded functions $\ell_q,r_q$ there, respectively lsc
and usc, such that eventually on the supports inside $U_q$,
\[
 \ell_q\le a_i\le r_q.
\]
Suppose $\ell_q(v),r_q(v)$ tend to the same number $a(v)$ for every
$v\in S$. Then
\[
                   \int a_i\,d\nu_i\longrightarrow\int_S a\,d\nu.
 \tag{2}
\]
The neighborhoods and the eventual index may depend on $q$.
\end{lemma}

\paragraph{Proof.}
Fix $q$. By normality of compact Hausdorff spaces and the finiteness of
$S$, choose a continuous $0\le\chi_q\le1$, equal to one on $S$, with
support in $U_q$. Clipping both $\ell_q$ and $r_q$ to $[-C,C]$
preserves their respective semicontinuity and their inequalities on
the supports. It also preserves their endpoint limits, since
$a(v)\in[-C,C]$ for $v\in S$.
The functions $\chi_q\ell_q,\chi_qr_q$, extended by zero outside
$U_q$, are bounded lsc and usc: at the boundary their absolute values
are bounded by $C\chi_q$, which tends to zero. The Radon weak-convergence
inequalities on compact spaces give
\[
 \int\ell_q\,d\nu
 \le\liminf_i\int\chi_q a_i\,d\nu_i
 \le\limsup_i\int\chi_q a_i\,d\nu_i
 \le\int r_q\,d\nu.
 \tag{3}
\]
These inequalities hold for nets as well. For example, a bounded lsc
function is the supremum of continuous minorants; approximate its Radon
integral from below by finitely many such minorants and use their maximum.
The usc inequality follows by changing sign. Thus no sequential
metrizability assumption is hidden.

Since $\nu$ is supported where $\chi_q=1$,
\[
 \left|\int(1-\chi_q)a_i\,d\nu_i\right|
        \le C\left(\nu_i(K)-\int\chi_q\,d\nu_i\right)\to0.
\]
Apply (3), then let $q$ tend along its directed degrees. The finitely
many endpoint values squeeze the integral to (2). \qed

\section{Apply the squeeze to the changing complete-system quotient}
Fix two rational model obstacles $f_D,f_H$, and first normalize their
envelopes to generic value zero. Ordinary constant shifts are harmless;
if their normalization constant is real, perform rational shifts first
and pass to the constant-equivariant formulas. No claim of an actual
rational line bundle with an irrational Cartier coefficient is needed.

Choose $C$ bounding $|f_H|$. Comparison of actual lattices gives
\[
       \phi_{m,0}-C\le\phi_{m,H}\le\phi_{m,0}+C.
 \tag{4}
\]
Indeed multiplying the zero-obstacle lattice by an appropriate time
power gives a sublattice of the $H$-obstacle lattice, and conversely;
the full section metric preserves these inequalities. Denominator
clearing is used in the degree. After normalization enlarge $C$ by the
absolute value of the constant shift. At auxiliary divisorial points
(1) therefore gives $|\psi_{m,H}|\le C$.

For every divisorial $v$, the nonnegligibility application in foundation
2 and foundation 1 give
\[
                \phi_{q,H}(v)\uparrow u_H(v),\qquad
                \phi_{q,0}(v)\uparrow V_\theta(v).
 \tag{5}
\]
Both limiting values are finite. Every nonzero complete-system section
has finite valuation at a function-field divisorial point, and one
sufficiently large nonempty degree gives a finite lower bound. Upper
bounds follow from the generic normalization and the genuine cone.
This is not a claim of equality of the unregularized supremum and its
usc regularization at supported points.

Fix a nonempty $q$. Full section metrics are continuous extended-valued
maxima on $K$, and finite in some ordinary neighborhood $U_q$ of each
of the finitely many output points $v_j$. Choose a common neighborhood
of them where both $\phi_{q,H}$ and $\phi_{q,0}$ are finite. On this
neighborhood both $u_H,V_\theta$ are finite too, since they dominate
those finite stages. Their upper bounds and the continuous finite lower
bounds permit shrinking around $S=\{v_j:c_j>0\}$ to obtain bounded
functions below. Define there
\[
 \ell_q=\max(-C,\phi_{q,H}-V_\theta),\qquad
 r_q=\min(C,u_H-\phi_{q,0}).                           \tag{6}
\]
The first is lsc and the second usc. For every sufficiently divisible
$m$ with $q\mid m$, monotonicity and the envelope majorizations imply
\[
 \ell_q\le \phi_{m,H}-\phi_{m,0}\le r_q\quad\hbox{on }U_q.
 \tag{7}
\]
There is no subtraction of two infinities on $U_q$.

The original valuation underlying every atom of $\mu_m$ has a unique
valuative lift on the birational generic model. The horizontal factor
has finite valuation there, so (1) identifies its quotient in (7) with
$\psi_{m,H}$ at that lift. Thus it defines a bounded measurable $a_m$
on the finite support of $\mu_m$. Arbitrary values elsewhere are not
used. By (5)--(6), the endpoint bounds at $S$ tend to
$u_H(v)-V_\theta(v)$; (4) keeps that difference in $[-C,C]$.
Lemma~\ref{squeeze} gives
\[
 \lim_m\int\psi_{m,H}\,d\MA_{A_m,\rm raw}(\psi_{m,D})
       =\sum_{j=1}^N c_j\bigl(u_H(v_j)-V_\theta(v_j)\bigr).
 \tag{8}
\]
This is an original-carrier weighted identification for these fixed
atomic outputs. It uses fixed-degree localized semicontinuous tests,
not an unsupported assertion of uniform quotient convergence.

\section{Identify the algebraic moment and extend its test}
For one sufficiently large common time twist put
$\alpha_Z=\pi_X^*L+Z+cF$, $Z=0,D,H$. The separately reconstructed
simultaneous mixed-moving limit is
\[
 \lim_m\int(-\psi_{m,H})\,d\MA_{A_m,\rm raw}(\psi_{m,D})
  =\langle\alpha_0\alpha_D^n\rangle
                        -\langle\alpha_H\alpha_D^n\rangle.
 \tag{9}
\]
These are mixed ordinary positive products; their difference's time-twist
independence comes from the actual finite moments. Under foundations
1--4, (8) proves exactly
\[
 M_{\rm alg}(H,D)=\int(V_\theta-u_H)\,d\mu_D.
 \tag{10}
\]
The reference-pole set is $\mu_D$-null, since it contains none of the
finitely many divisorial atoms where (5) gives a finite value.

For normalized bounded nef model potentials on generic $A$, the
separate finite Hodge estimate is
\[
 \int(-h)\,d\MA_{A,\rm raw}(\psi)
 \le(-g_h)A^n+C_nJ_A(h)^{\alpha_n}J_A(\psi)^{1/2}
                  \max\{J_A(h),J_A(\psi)\}^{1/2-\alpha_n},
 \quad\alpha_n=2^{-n},
 \tag{11}
\]
where $g_h=h(v_0)\le0$ and $J_A(h)=g_hA^n-e_A(h)$.
All masses are raw. Taking the moving limit before the genuine test
limit gives, for a normalized model-envelope test $u_H$,
\[
 \int(V_\theta-u_H)\,d\mu_D
 \le C_n B(u_H)^{\alpha_n}B(u_D)^{1/2}
                   \max\{B(u_H),B(u_D)\}^{1/2-\alpha_n}.
 \tag{12}
\]
Here $B(u)=-E(u)$ and the generic error in (11) tends to zero.

Now take any genuine normalized $v\in\mathcal E^1_\theta$, and choose
model envelopes $v_j\downarrow v$ with $E(v_j)\downarrow E(v)$.
Put $d_j=v_j(v_0)\downarrow0$. Apply (12) to $v_j-d_j$; these need
not decrease. Their costs $-E(v_j)+d_jV$ are bounded and tend to $B(v)$.
On the finite set $S$ the unnormalized $v_j$ decrease. Therefore
$V_\theta-v_j$ increases; adding one fixed constant makes all these
functions nonnegative on $S$. Ordinary monotone integration and
$d_j\to0$ imply
\[
 \int(V_\theta-v)\,d\mu_D
 \le C_n B(v)^{\alpha_n}B(u_D)^{1/2}
                    \max\{B(v),B(u_D)\}^{1/2-\alpha_n}<\infty.
 \tag{13}
\]
Consequently every such $v$ is finite at every atom with $c_j>0$.
Finite-energy tests for a fixed finite atomic output have thus been
identified and controlled under the explicit foundations. No monotonicity
of the normalized envelopes or of the changing section quotient was used.

\paragraph{Exact residual boundary.}
The proof depends essentially on a fixed finite divisorial support:
it puts $\chi_q=1$ on the entire limit measure, and takes endpoint limits
at finitely many finite values. It does not establish pole-nullness of a
weak limit of outputs of arbitrary unbounded potentials. Nor does it
prove weighted convergence for such non-atomic outputs, a uniform
absolute continuous-test approximation, strong integration against an
arbitrary finite-energy target, strict comparison or an actual attained
inverse equation. These are still required by original P16.5.
The foundations 1--4 and the new local squeeze application require
independent verification. This supporting conditional paper closes no
original problem and changes no accepted-paper count.

\begin{thebibliography}{9}
\bibitem{BJ} S. Boucksom and M. Jonsson,
\emph{Global pluripotential theory over a trivially valued field},
\href{https://arxiv.org/pdf/1801.08229v4}{arXiv:1801.08229v4},
Section 5.1 and Theorem 5.6: divisorial nonnegligibility. Their ample
application must be reconstructed before foundation 2 is used.
\bibitem{bridges} Full relative-floor closure reconstruction;
complete section-volume/model-output reconstruction; finite-nef moment
and mixed-moving-limit reconstruction, October 2026. Their exact
mathematical statements are explicit conditional foundations here,
not conclusions supplied by a historical verdict or a file hash.
\end{thebibliography}
\end{document}
